\begin{statementbox}
	\textbf{\textcolor{CStatement}{条件}}
	\begin{description}[style=nextline, leftmargin=2.8em, labelsep=0.5em, itemsep=0.35em]
		\item[\textbf{\textcolor{CStatement}{条件1（奇偶性与定义域）：}}] 函数 $f(x)$ 为奇函数或偶函数，且定义域关于原点对称。
		\item[\textbf{\textcolor{CStatement}{条件2（半轴单调性）：}}] 在区间 $[0,+\infty)$ 上 $f(x)$ 严格单调递增。
	\end{description}

	\textbf{\textcolor{CStatement}{结论}}
	\begin{description}[style=nextline, leftmargin=2.8em, labelsep=0.5em, itemsep=0.45em]
		\item[\textbf{\textcolor{CStatement}{结论一（奇函数等价）：}}]
			\textbf{\textcolor{CStatement}{直观描述：}}
			奇函数在正半轴递增，则函数值之和的正负与自变量之和一致。
			\[
				f(a)+f(b)>0 \iff a+b>0 .
			\]

		\item[\textbf{\textcolor{CStatement}{结论二（偶函数等价）：}}]
			\textbf{\textcolor{CStatement}{直观描述：}}
			偶函数在正半轴递增，则函数值大小与自变量绝对值大小一致。
			\[
				f(a)>f(b) \iff |a|>|b| .
			\]
	\end{description}

	\textbf{\textcolor{CStatement}{推广结论（单调递减反向）：}}
	\textbf{\textcolor{CStatement}{直观描述：}}
	若 $[0,+\infty)$ 上单调递减，则奇函数结论变为和小于零，偶函数结论变为绝对值小于关系。
	\[
		f(a)+f(b)>0 \iff a+b<0,
	\]
	\[
		f(a)>f(b) \iff |a|<|b|.
	\]
\end{statementbox}
